\scanpage{331}
\mspage{3}

dans \deleted{$\operatorname{Hom}_{\Pi(X)}$} $\mathrm{Fl}\,\Pi(X)$ est réduite à un seul élément.

Ainsi on trouve que \uline{$X$ est 1-connexe}. Il y a donc de l'espoir ! Il reste à prouver

\[
 H^i(X,M)=0 \qquad i\geq 2,\quad M\ \text{coeff. const.}
\]

J'\deleted{\unsure{y}}ai \deleted{vérifié lorsque $b=c$, et \unsure{on}} \inserted{regardé le cas lorsqu'on} a ceci : $b=c$, $\mathrm{Ob}\,X=\{b,b\times b\}$, $\operatorname{Hom}_X(b\times b)=\{\text{réunion disjointe des ens }u\circ\mathrm{pr}_1\text{ et }u\circ\mathrm{pr}_2,\text{ où }u\in\operatorname{Hom}_X(b,b)=G,\text{ un monoïde quelconque.}$ \deleted{\illegible{1}} \deleted{Alors} \inserted{\unsure{Ici}} $X^{\unsure{\circ|o}}$ est \deleted{\illegible{1}} fibrée, \inserted{via $I\mapsto b^I$,} sur la \inserted{sous} catégorie pleine \inserted{$S$} de Ens formée de $\{0\}$ et $\{0,1\}$, et

\[
 \deleted{H^n(X,M)\simeq H^0(S,\mathscr H^n)\quad\unsure{\text{or ici}}}
\]

\deleted{et on vérifie \inserted{surtout} que c'est nul pour $n\geq 1$. (NB. les $H^p(S,\unsure{\mathscr H^q})$ sont nuls pour $p>0$, car $S$ a $\{0\}$ comme objet final. \unsure{Si}...} Donc $X$ est cofibrée (\deleted{\unsure{sur}} \unsure{potentiellement|particules|parties} propres) sur $S^\circ\overset{\text{déf}}{=}T$. (NB. $\unsure{S^\circ|S^0}=T$ a un objet initial, mais non un objet final. Sa \unsure{coh.|cat.} à val. dans un préfaisceau quelc. \unsure{n'est} pas \unsure{tjs} \deleted{\unsure{bien connu}} triviale, cf ci dessous)

\[
 H^n(X,M)\Longleftarrow H^p\bigl(T,\underbrace{R^q f_*(M_X)}_{=\mathscr H^q}\bigr)
 \qquad \text{où }f:X\unsure{\text{---}|\longrightarrow} S^\circ .
\]

J'admet ici (provisoirement) ceci :

\[
 \begin{aligned}
 \deleted{\unsure{\forall}}\ H^p(\deleted{\unsure{S^\circ}}T,\deleted{\unsure{\mathscr H^q|R^q}}F)&=0\qquad\text{pour }p\geq2,\ \text{tout }F\in T^\wedge.\\
 H^0(T,F)&\simeq\operatorname{Ker}\delta\\
 H^1(T,F)&\simeq\operatorname{Coker}\delta
 \end{aligned}
\]

\margin{où $\delta:F_0\longrightarrow F_1$ est défini comme la différence $\delta=p_1^*-p_2^{\unsure{*|-}}$}

\margin{où $p_1,p_2$ sont les deux proj. de $\{0,1\}\simeq\{0\}^2$ dans $T^{\unsure{\circ}}\!=S^\circ$}

\deleted{On vérifie de plus} \inserted{On a de plus}

\[
 (\mathscr H^q)_0=H^q(G,M),\quad (\mathscr H^q)_1=H^q(G\times G,M)
\]

et

\[
 \mathrm{pr}_1^*,\mathrm{pr}_2^*:H^q(G,M)\rightrightarrows H^q(\unsure{G\times G},M)
\]

sont les \unsure{transposés|transportés|transposées} de $\mathrm{pr}_1,\mathrm{pr}_2:G\times G\longrightarrow G$

% archive sheet number: 331
